% Part: incompleteness
% Chapter: representability-in-q
% Section: basic-representable

\documentclass[../../../include/open-logic-section]{subfiles}

\begin{document}

% SEGMENT OLP-0294-S01

\olfileid{inc}{req}{bre}
\olsection{அடிப்படைச் சார்புகள்~$\Th{Q}$ இல் பிரதிநிதித்துவப்படுத்தத்தக்கவை}

முதலில் எல்லா அடிப்படைச் சார்புகளும்~$\Th{Q}$ இல்
பிரதிநிதித்துவப்படுத்தத்தக்கவை என்று காட்ட வேண்டும். இறுதியில், ஒவ்வொரு
$k$-இட அடிப்படைச் சார்பு $f(x_0,\dots,x_{k-1})$ க்கும், அதைப்
பிரதிநிதித்துவப்படுத்தும் !!a{formula}~$!A_f(x_0,\dots,x_{k-1},y)$
ஒன்றை எவ்வாறு ஒதுக்குவது என்று காட்ட வேண்டும்.

சுழியம், அடுத்தெண், கூட்டல், பெருக்கல், சமத்துவத்தின் சிறப்பியல்புச் சார்பு,
வீழ்த்தல்கள் ஆகியவற்றை நம்மால் பிரதிநிதித்துவப்படுத்த முடியும். ஒவ்வொரு
நேர்விலும் பொருத்தமான பிரதிநிதித்துவ வாய்பாடு முற்றிலும் நேரடியானது;
எடுத்துக்காட்டாக, சுழியம் $y = \Obj 0$ என்ற வாய்பாட்டாலும், அடுத்தெண்
$x_0' = y$ என்ற !!{formula} வாலும், கூட்டல் $(x_0 + x_1) = y$ என்ற
!!{formula} வாலும் பிரதிநிதித்துவப்படுத்தப்படுகின்றன. தொடர்புடைய
!!{sentence}s ஐ~$\Th{Q}$ நிறுவ முடியும் என்று காட்டுவதே இதில் உள்ள பணி;
எடுத்துக்காட்டாக, மேலுள்ள !!{formula} ஆனது கூட்டலைப்
பிரதிநிதித்துவப்படுத்துகிறது என்று கூறுவதில், ஒவ்வோர் இயல் எண் சோடி
$m$,~$n$ க்கும் பின்வருவனவற்றை $\Th{Q}$ நிறுவுகிறது என்று காட்டுவதும்
அடங்கும்:
\begin{align*}
& \eq[\num n + \num m][\num {n+m}] \text{ மற்றும்}\\
& \lforall[y][(\eq[(\num n + \num m)][y] \lif \eq[y][\num{n+m}])].
\end{align*}

\begin{prop}
\ollabel{prop:rep-zero}
சுழியச் சார்பு $\Zero(x) = 0$ ஆனது~$\Th{Q}$ இல்
$!A_{\Zero}(x,y) \ident \eq[y][\Obj 0]$ என்பதால் பிரதிநிதித்துவப்படுத்தப்படுகிறது.
\end{prop}

\begin{prop}
\ollabel{prop:rep-succ}
அடுத்தெண் சார்பு $\Succ(x) = x+1$ ஆனது~$\Th{Q}$ இல்
$!A_{\Succ}(x,y) \ident \eq[y][x']$ என்பதால் பிரதிநிதித்துவப்படுத்தப்படுகிறது.
\end{prop}

\begin{prop}
\ollabel{prop:rep-proj}
வீழ்த்தல் சார்பு $\Proj{n}{i}(x_0, \dots, x_{n-1}) = x_i$ ஆனது
~$\Th{Q}$ இல்
\[!A_{\Proj{n}{i}}(x_0, \dots, x_{n-1}, y) \ident \eq[y][x_i].\]
என்பதால் பிரதிநிதித்துவப்படுத்தப்படுகிறது.
\end{prop}

\begin{prob}
$\eq[y][\Obj 0]$, $\eq[y][x']$, $\eq[y][x_i]$ ஆகியவை முறையே
$\Zero$, $\Succ$, $\Proj{n}{i}$ ஆகியவற்றைப் பிரதிநிதித்துவப்படுத்துகின்றன
என்று நிறுவுக.
\end{prob}

\begin{prop}
\ollabel{prop:rep-id}
~$=$ இன் சிறப்பியல்புச் சார்பு,
% SOURCE CORRECTION TA-INC-019: localize the bare math-mode prose in the second case.
\[
\Char{=}(x_0, x_1) =
\begin{cases}
  1 & x_0 =x_1 \text{ எனில்}\\
  0 & \text{இல்லையெனில்}
\end{cases}
\]
ஆனது~$\Th{Q}$ இல் பின்வருவதால் பிரதிநிதித்துவப்படுத்தப்படுகிறது:
\[
  !A_{\Char{=}}(x_0, x_1, y) \ident (\eq[x_0][x_1] \land \eq[y][\num{1}]) \lor (\eq/[x_0][x_1] \land
\eq[y][\num{0}]).
\]
\end{prop}

இதன் நிறுவலுக்குப் பின்வரும் துணைத்தேற்றம் தேவைப்படுகிறது.

\begin{lem}
\ollabel{lem:q-proves-neq} இயல் எண்கள் $n$,~$m$ தரப்பட்டிருக்கும்போது,
$n \neq m$ எனில், $\Th{Q} \Proves \eq/[\num n][\num m]$.
\end{lem}

\begin{proof}
ஒவ்வொரு $m$ க்கும், $n \neq m$ எனில், $Q \Proves
\eq/[\num n][\num m]$ என்று காட்ட $n$ மீது தொகுத்தறிதலைப் பயன்படுத்துக.

அடிப்படை நேர்வில், $n = 0$. $m$ ஆனது $0$ க்குச் சமமில்லை எனில்,
ஏதாவது இயல் எண் $k$ க்கு $m = k + 1$. $\lforall[x][\eq/[0][x']]$
என்று கூறும் அடிகோள் நம்மிடம் உள்ளது. ஓர் அளவையடை அடிகோளால் $x$ க்குப்
பதிலாக $\num k$ ஐ இட்டு, $\eq/[0][\num k']$ என்று முடிவுசெய்யலாம்.
ஆனால் $\num k'$ என்பது வெறுமனே $\num m$.

தொகுத்தறிதல் படியில், கூற்று $n$ க்கு மெய் என்று எடுகொண்டு, $n+1$ ஐக்
கருதலாம். $m$ ஏதேனும் இயல் எண் ஆகட்டும். இரண்டு வாய்ப்புகள் உள்ளன:
$m = 0$ அல்லது ஏதாவது $k$ க்கு $m = k+1$. முதல் நேர்வு மேலுள்ளவாறே
கையாளப்படுகிறது. இரண்டாவது நேர்வில், $n+1 \neq k+1$ என்று கொள்க.
அப்போது $n \neq k$. $n$ க்கான தொகுத்தறிதல் எடுகோளின்படி,
$\Th{Q} \Proves \eq/[\num n][\num k]$. நம்மிடம்
$\lforall[x][\lforall[y][\eq[x'][y'] \lif \eq[x][y]]]$ என்று கூறும்
ஓர் அடிகோள் உள்ளது. ஓர் அளவையடை அடிகோளைப் பயன்படுத்தி,
$\eq[\num n'][\num k'] \lif \eq[\num n][\num k]$ ஐப் பெறுகிறோம்.
கூற்றுத் தருக்கத்தைப் பயன்படுத்தி, $\Th{Q}$ இல்
$\eq/[\num n][\num k] \lif \eq/[\num n'][\num k']$ என்று முடிவுசெய்யலாம்.
Modus ponens ஐப் பயன்படுத்தி, நாம் விரும்பிய $\eq/[\num n'][\num k']$
ஐ முடிவுசெய்யலாம்; ஏனெனில் $\num k'$ என்பது $\num m$.
\end{proof}

\begin{explain}
இத்துணைத்தேற்றம் அதிகம் கூறவில்லை என்பதைக் கவனிக்க: வெவ்வேறு எண்ணுருக்கள்
வெவ்வேறு பொருள்களைக் குறிக்கின்றன என்பதை~$\Th{Q}$ நிறுவ முடியும் என்றே
அது சாரத்தில் கூறுகிறது. எடுத்துக்காட்டாக, $0'' \neq 0'''$ என்பதை
$\Th{Q}$ நிறுவுகிறது. ஆனால் இது பொதுவாகப் பொருந்துகிறது என்று காட்டுவதில்
சிறிது கவனம் தேவை. நாம் தொகுத்தறிதலைப் பயன்படுத்தினாலும், அது $\Th{Q}$
க்கு \emph{வெளியே} மேற்கொள்ளப்படும் தொகுத்தறிதல் என்பதையும் கவனிக்க.
\end{explain}

\begin{proof}[\olref{prop:rep-id} இன் நிறுவல்]
$n = m$ எனில், $\num{n}$,~$\num{m}$ ஆகியவை ஒரே சொல்லாகும்; மேலும்
$\Char{=}(n, m) = 1$. ஆனால் $\Th{Q} \Proves
(\eq[\num{n}][\num{m}] \land \eq[\num{1}][\num{1}])$; ஆகவே அது
$!A_=(\num{n}, \num{m}, \num{1})$ ஐ நிறுவுகிறது. $n \neq m$ எனில்,
$\Char=(n, m) = 0$. \olref{lem:q-proves-neq} இன்படி,
$\Th{Q} \Proves \eq/[\num{n}][\num{m}]$; எனவே
$(\eq/[\num{n}][\num{m}] \land \Obj 0 = \Obj 0)$ ஐயும் நிறுவுகிறது.
ஆகவே $\Th{Q} \Proves !A_=(\num{n}, \num{m}, \num{0})$.

இரண்டாம் பகுதிக்கும் இரண்டு நேர்வுகள் உள்ளன. $n = m$ எனில்,
$\Th{Q} \Proves \lforall[y][(!A_=(\num{n}, \num{m}, y)
  \lif \eq[y][\num{1}])]$ என்று காட்ட வேண்டும். முறைசாரா வகையில்
வாதிடுகையில், $!A_=(\num{n}, \num{m}, y)$, அதாவது பின்வருவது என்று கொள்க:
\[
(\eq[\num{n}][\num{n}] \land \eq[y][\num{1}]) \lor
(\eq/[\num{n}][\num{n}] \land \eq[y][\num{0}])
\]
இடப்பக்கப் பிரிப்புக்கூறு தருக்கத்தின்படி $\eq[y][\num{1}]$ ஐப்
பின்விளைவிக்கிறது; வலப்பக்கப் பிரிப்புக்கூறு தருக்கத்தால் நிறுவத்தக்க
$\eq[\num{n}][\num{n}]$ உடன் முரண்படுகிறது.

மறுபுறம், $n \neq m$ என்று கொள்க. அப்போது $!A_=(\num{n},
\num{m}, y)$ என்பது
\[
(\eq[\num{n}][\num{m}] \land \eq[y][\num{1}]) \lor
(\eq/[\num{n}][\num{m}] \land \eq[y][\num{0}])
\]
ஆகும். இங்கு இடப்பக்கப் பிரிப்புக்கூறு, \olref{lem:q-proves-neq} இன்படி
$\Th{Q}$ இல் நிறுவத்தக்க $\eq/[\num{n}][\num{m}]$ உடன் முரண்படுகிறது;
வலப்பக்கப் பிரிப்புக்கூறு $\eq[y][\num{0}]$ ஐப் பின்விளைவிக்கிறது.
\end{proof}

\begin{prop}
\ollabel{prop:rep-add}
கூட்டல் சார்பு $\Add(x_0, x_1) = x_0+x_1$ ஆனது~$\Th{Q}$ இல்
பின்வருவதால் பிரதிநிதித்துவப்படுத்தப்படுகிறது:
\[
  !A_{\Add}(x_0, x_1, y) \ident \eq[y][(x_0 + x_1)].
\]
\end{prop}

\begin{lem}
\ollabel{lem:q-proves-add}
$\Th{Q} \Proves \eq[(\num{n} + \num{m})][\num{n+m}]$
\end{lem}

\begin{proof}
$m$ மீது தொகுத்தறிதலால் இதை நிறுவுகிறோம். $m = 0$ எனும் நேர்வில்,
$\Th{Q} \Proves \eq[(\num{n} + \Obj 0)][\num{n}]$ என்பதே கூற்று.
இது அடிகோள்~$!Q_4$ இன்படி பின்வருகிறது. இப்போது $m$ க்கான கூற்றை
எடுகொண்டு, $m+1$ க்கான கூற்றை நிறுவுவோம்; அதாவது,
$\Th{Q} \Proves \eq[(\num{n} + \num{m+1})][\num{n+m+1}]$ என்று
நிறுவுவோம். $\num{m+1}$ என்பது வெறுமனே $\num{m}'$ என்றும்,
$\num{n+m+1}$ என்பது வெறுமனே $\num{n+m}'$ என்றும் கவனிக்க. அடிகோள்
$!Q_5$ இன்படி, $\Th{Q} \Proves \eq[(\num{n} +
\num{m}')][(\num{n}+\num{m})']$. தொகுத்தறிதல் எடுகோளின்படி,
$\Th{Q} \Proves \eq[(\num{n} + \num{m})][\num{n+m}]$. ஆகவே
$\Th{Q} \Proves \eq[(\num{n} + \num{m}')][\num{n+m}']$.
\end{proof}

\begin{proof}[\olref{prop:rep-add} இன் நிறுவல்]
$\Add$ ஐப் பிரதிநிதித்துவப்படுத்தும் !!{formula}~$!A_\Add(x_0, x_1, y)$
என்பது $\eq[y][(x_0 + x_1)]$. முதலில் $\Add(n, m) = k$ எனில்,
$\Th{Q} \Proves !A_\Add(\num{n}, \num{m}, \num{k})$, அதாவது
$\Th{Q} \Proves \eq[\num{k}][(\num{n} + \num{m})]$ என்று காட்டுகிறோம்.
ஆனால் $k = n + m$ என்பதால், $\num{k}$ என்பது $\num{n+m}$ தான்;
மேலும் \olref{lem:q-proves-add} இல் $\Th{Q} \Proves
\eq[(\num{n} + \num{m})][\num{n+m}]$ என்று காட்டியுள்ளோம்.

$\Add(n, m) = k$ எனில், பின்வருவதையும் காட்ட வேண்டும்:
\[
\Th{Q} \Proves \lforall[y][(!A_\Add(\num{n}, \num{m}, y) \lif
  \eq[y][\num{k}])].
\]
$\eq[(\num{n} + \num{m})][y]$ நம்மிடம் உள்ளது என்று கொள்க.
பின்வருவது உள்ளதால்,
\[
\Th{Q} \Proves \eq[(\num{n}+\num{m})][\num{n+m}],
\]
இடப்பக்கத்திற்குப் பதிலாக $\num{n+m}$ ஐ இட்டு, தன்னிச்சையான~$y$ க்கு
$\eq[\num{n+m}][y]$ என்று பெறலாம்.
\end{proof}

\begin{prop}
\ollabel{prop:rep-mult}
பெருக்கல் சார்பு $\Mult(x_0, x_1) = x_0 \cdot x_1$ ஆனது~$\Th{Q}$ இல்
பின்வருவதால் பிரதிநிதித்துவப்படுத்தப்படுகிறது:
\[
  !A_{\Mult}(x_0, x_1, y) \ident y = (x_0 \times x_1).
\]
\end{prop}

\begin{proof} பயிற்சி. \end{proof}

\begin{lem}
\ollabel{lem:q-proves-mult}
$\Th{Q} \Proves \eq[(\num{n} \times \num{m})][\num{n \cdot m}]$
\end{lem}

\begin{proof} பயிற்சி. \end{proof}

\begin{prob}
\olref[inc][req][bre]{lem:q-proves-mult} ஐ நிறுவுக.
\end{prob}

\begin{prob}
\olref[inc][req][bre]{lem:q-proves-mult} ஐப் பயன்படுத்தி,
\olref[inc][req][bre]{prop:rep-mult} ஐ நிறுவுக.
\end{prob}

\begin{explain}
  எண்கணித மொழியின் சார்புக் குறிக்கு $\times$ ஐயும், எண்களின் மீதான
  சாதாரணப் பெருக்கல் செயலுக்கு $\cdot$ ஐயும் பயன்படுத்துகிறோம் என்பதை
  நினைவுகூர்க. எனவே எண்களைக் குறிக்கும் கோவைகளுக்கு இடையில் (எடுத்துக்காட்டாக,
  $m \cdot n$ இல்) $\cdot$ இடம்பெறலாம்; ஆனால் எண்கணித மொழியின் சொற்களுக்கு
  இடையில் மட்டும் (எடுத்துக்காட்டாக, $(\num{m} \times \num{n})$ இல்)
  $\times$ இடம்பெறும். இன்னும் குழப்பமூட்டும் வகையில், !!{function}
  மற்றும் கூட்டல் செயல் இரண்டுக்கும் $+$ பயன்படுத்தப்படுகிறது. சொற்களுக்கு
  இடையில்---எடுத்துக்காட்டாக, $(\num{n} + \num{m})$ இல்---அது எண்கணித
  மொழியின் $2$-இட !!{function}; எண்களுக்கு இடையில்---எடுத்துக்காட்டாக,
  $n+m$ இல்---அது கூட்டல் செயல். இதில் $\num{n+m}$ என்ற நேர்வும் அடங்கும்:
  இது எண்~$n+m$ க்கு ஒத்த தரநிலை எண்ணுரு.
\end{explain}

\end{document}
